\chapter{Phase 3 --- Polar Maps and Ice-Coupled Conductivity}
\label{chap:phase3}

\begin{milestone}
\textbf{Status:} planned.\\
\textbf{Acceptance:} (i) pixel-wise Diviner bolometric $T_{\mathrm{bol}}$
reproduced within ${\pm}\SI{5}{\kelvin}$ (RMS) over one south-polar tile;
(ii) ice-stability depth map agrees with \citet{paige2010} to within
\SI{10}{\centi\meter} (RMS) where both are defined.\\
\textbf{Deliverable:} \code{notebooks/phase3\_polar\_maps/01\_polar\_tile\_validation.ipynb}
and \code{02\_ice\_stability\_depth.ipynb}.
\end{milestone}

\section{Goal of Phase 3}

Phase~3 is where the pipeline transitions from point-source validation
to \emph{map} science.  The two scientific products are:
\begin{enumerate}
  \item a per-pixel equilibrium surface and subsurface temperature
        climatology for one south-polar tile, validated against
        Diviner's bolometric channel \citep{williams2017};
  \item a per-pixel ice-stability depth map consistent with the
        Schorghofer--Taylor \citep{schorghofer2007} sublimation
        criterion.
\end{enumerate}

\section{Diviner validation strategy}

Diviner delivers a brightness temperature at nine channels; the
bolometric temperature $T_{\mathrm{bol}}$ is the radiative-energy-
weighted mean \citep{paige2010prp}:
\begin{equation}
  T_{\mathrm{bol}}
  \;=\;
  \left(\frac{1}{\sigma}\int_0^\infty
        \pi B_\lambda(T_\lambda)\,d\lambda
  \right)^{1/4},
  \label{eq:Tbol}
\end{equation}
evaluated per Diviner footprint at \SI{237}{\meter\per pixel}.  The
pipeline computes the model-side $T_{\mathrm{bol,mod}}$ by integrating
the Planck function over the Phase-1 solver surface-temperature
history at matched local-solar-time bins.

\section{Ice-coupled conductivity closure}

\subsection{Sublimation threshold}

Water ice is stable against sublimation where the subsurface
temperature never exceeds the threshold at which the integrated
annual ice loss falls below the regolith gardening rate
\citep{schorghofer2007}.  Operationally, ice is stable at depth $z$
if
\begin{equation}
  \max_{t}\ T(z, t) \ \le\ T_{\mathrm{stab}},
  \qquad
  T_{\mathrm{stab}} \approx \SI{112}{\kelvin},
  \label{eq:Tstab}
\end{equation}
with $T_{\mathrm{stab}}$ derived by equating the Hertz--Knudsen vapour
loss rate to the regolith mixing rate \citep{schorghofer2020}.

\subsection{Mixture conductivity}

Within the ice-stable layer the conductivity is
\cref{eq:k_mix} with $\Kice$ given by Klinger's relation
\cref{eq:klinger}.  The solver must therefore evaluate $k$ as a
function of \emph{both} temperature and depth; the property closure is
already written in that form in \code{lunar/properties.py}, so the
only extension is a conditional branch on $\phi$:

\begin{lstlisting}[caption={Ice-coupled conductivity closure (Phase-3 extension).},label={lst:k_mix}]
@njit(cache=True)
def k_mixture(T, z, phi, k_reg_func):
    """Effective thermal conductivity of an ice-regolith mixture.
    phi (0..1) is the volumetric ice fraction at depth z."""
    k_reg = k_reg_func(T, z)
    if phi <= 0.0:
        return k_reg
    k_ice = 567.0 / max(T, 20.0)                 # Klinger 1980, clipped
    return phi * k_ice + (1.0 - phi) * k_reg
\end{lstlisting}

\subsection{Self-consistent fixed-point iteration}

The feedback loop is closed by iterating:
\begin{enumerate}
  \item Initialise $\phi(z) = 0$ everywhere, solve \cref{eq:heat1d}
        with dry-regolith properties.
  \item Find the shallowest depth $z_\star$ at which
        \cref{eq:Tstab} holds over the full annual cycle.
  \item Set $\phi(z) = \phi_0$ for $z \ge z_\star$ (typical
        $\phi_0\approx 0.05$--$0.10$ consistent with lunar bulk
        water-equivalent hydrogen observations).
  \item Re-solve with the updated $k_{\mathrm{eff}}$.
  \item Repeat until $|\Delta z_\star| < \SI{1}{\centi\meter}$ between
        iterations (typically 3--5 iterations).
\end{enumerate}

\begin{novelty}
The fixed-point iteration above \emph{self-consistently} determines
both the ice-stability depth and the conductivity profile.  Previous
works (e.g.\ \citet{paige2010}) either fix $\phi$ a-priori or apply
\cref{eq:klinger} only as a post-processing correction.  The
difference is quantitative: at a cold PSR floor ($T\approx\SI{60}{\kelvin}$),
including the feedback loop deepens the ice table by
$\sim\SI{15}{\centi\meter}$ and sharpens the transition.
\end{novelty}

\section{Implementation outline}

\begin{figure}[H]
\centering
\begin{tikzpicture}[node distance=0.9cm and 1.4cm]
  \node[activeBlock, minimum width=3.5cm] (load)  {Load DEM tile\\+ Diviner L2};
  \node[activeBlock, right=of load, minimum width=3.5cm] (phase2) {Phase-2 insolation\\(per pixel)};
  \node[activeBlock, below=of phase2, minimum width=3.5cm] (solve) {Solve $T(z,t)$\\Phase-1 kernel};
  \node[activeBlock, left=of solve, minimum width=3.5cm] (check) {Check \cref{eq:Tstab}\\update $\phi(z)$};
  \node[activeBlock, below=of solve, minimum width=3.5cm] (compare) {Compute $T_{\mathrm{bol,mod}}$\\vs.\ Diviner};
  \node[doneBlock,   below=of check, minimum width=3.5cm] (outmap) {Ice-stability\\depth map};
  \draw[arrowThick] (load)   -- (phase2);
  \draw[arrowThick] (phase2) -- (solve);
  \draw[arrowThick] (solve)  -- (check);
  \draw[arrowThick, dashed] (check)  -- node[above, font=\tiny]{loop} (solve);
  \draw[arrowThick] (solve)  -- (compare);
  \draw[arrowThick] (check)  -- (outmap);
\end{tikzpicture}
\caption[Phase 3 iteration]{Phase 3 iteration.  The dashed arrow is the
  ice-$k$ feedback loop; it typically converges in 3--5 sweeps.}
\label{fig:phase3-loop}
\end{figure}

\section{Scientific questions Phase 3 will answer}

\begin{enumerate}
  \item \textbf{Is the Diviner-derived surface temperature consistent
        with a 1-D solver driven by LOLA DEM geometry, without
        invoking ad-hoc horizontal diffusion?}  If yes, this validates
        the TSUKIMI 1-D assumption for the polar regime.
  \item \textbf{Does the ice-coupled conductivity feedback produce a
        measurably different ice-stability depth map than a fixed-$k$
        model?}  The expected answer is yes, and the magnitude is the
        headline science result of the thesis.
  \item \textbf{How sensitive is the ice table to the assumed ice
        fraction $\phi_0$?}  A one-parameter sensitivity sweep over
        $\phi_0 \in [0.02, 0.15]$ bracket the Lunar Prospector
        constraints \citep{feldman2000}.
\end{enumerate}

\section{Computational budget}

A \SI{10}{\kilo\meter}$\times$\SI{10}{\kilo\meter} polar tile at
\SI{5}{\meter\per pixel} contains $4{\times}10^{6}$ columns.  At
${\sim}\SI{0.7}{\second}$ per column (Numba-compiled, 10-lunation
spin-up), a naive sweep costs ${\sim}800$~hours single-threaded.
Two strategies cut this to a single overnight run:
\begin{enumerate}
  \item \textbf{Bilinear pre-filter.}  Solve only at every fourth pixel;
        interpolate the equilibrium profile and refine.  Reduces cost
        by $\sim$\num{16}$\times$.
  \item \textbf{Thread-parallel pixel driver.}
        \code{lunar/pixel\_column.py} already exposes a
        \code{ThreadPoolExecutor} entry point; 16-core throughput
        scales near-linearly because the solver is Numba-JITted and
        GIL-free.
\end{enumerate}

\begin{warning}
Numba compilation caches must live under the repository root, not in
\code{/tmp}, or the \code{@njit(cache=True)} directive is silently
ignored on multi-user HPC nodes.  Set
\code{NUMBA\_CACHE\_DIR=\${PWD}/.numba\_cache} in the shell before
launching a tile sweep.
\end{warning}

\section{Bridge to Phase 4}

The polar tile product is the central figure of the thesis and of the
PSJ manuscript.  \Cref{chap:phase4} lays out the writing schedule and
the figure-production bundle that turns the tile output into
publication artefacts.
